%% —-----------------------------------------------------------
%% sbcreviews-pt.tex
%% Template para SBC Reviews
%% (https://journals-sol.sbc.org.br/index.php/reviews/)
%%
%% Esta versão do template para a SBC Reviews inclui algumas
%% modificações sobre o template original da REIC
%% (https://journals-sol.sbc.org.br/index.php/reic)
%%
%% Licença CC BY 4.0 (https://creativecommons.org/licenses/by/4.0/)
%% Sociedade Brasileira de Computação.

\documentclass[portuguese]{sbcreviews-2025}%
\usepackage[misc,geometry]{ifsym}

\usepackage{tabularx}
\usepackage{aas_macros}
\usepackage[bottom]{footmisc}
\usepackage{graphicx}
\usepackage{tabularray}
\usepackage{afterpage}
\usepackage{url}
\usepackage{pifont}
\usepackage{listings}
\usepackage{xcolor}
\usepackage{booktabs}

\lstset{
    basicstyle=\ttfamily\scriptsize,
    breaklines=true,
    frame=single,
    columns=fullflexible
}
\setcitestyle{square}

\definecolor{engtitle}{rgb}{0.5,0.5,0.5}
\definecolor{orcidlogo}{rgb}{0.37,0.48,0.13}
\definecolor{unilogo}{rgb}{0.16, 0.26, 0.58}
\definecolor{maillogo}{rgb}{0.58, 0.16, 0.26}
\definecolor{darkblue}{rgb}{0.0,0.0,0.0}
\hypersetup{colorlinks,breaklinks,
            linkcolor=darkblue,urlcolor=darkblue,
            anchorcolor=darkblue,citecolor=darkblue}

\jid{PUC-Minas}
\jtitle{Trabalho Prático - Computação Distribuída, 2026}
\issn{2966-3938}
\doi{10.5753/rocs.2026.XXXXXX}
\copyrightstatement{This work is licensed under a Creative Commons Attribution 4.0 International License}
\jyear{2026}

\category{Sistemas Distribuídos}
\title[Voxel Royale Distribuído]{Backend Distribuído para um Jogo Multiplayer em Go e gRPC: Segunda Entrega do Voxel Royale}
\engtitle{\textcolor{engtitle}{A Distributed Backend for a Multiplayer Game in Go and gRPC: Second Delivery of Voxel Royale}}

%THE ORCID IS MANDATORY FOR EACH AUTHOR
\author[Silva et al. 2026]{
\affil{\textbf{Ana Cristina Martins Silva}~\orcidlink{0000-0000-0000-0000}~\textcolor{blue}{\faEnvelope}~~[{PUC Minas}~|\href{mailto: ana.silva.1353892@sga.pucminas.br}{~{\textit{ana.silva.1353892@sga.pucminas.br}}}~]}
\affil{\textbf{André Luiz Rocha Cabral}~\orcidlink{0000-0000-0000-0000}~[{PUC Minas}~|\href{mailto: alrcabral@sga.pucminas.br}{~{\textit{alrcabral@sga.pucminas.br}}}~]}
\affil{\textbf{João Paulo Dias Estevão}~\orcidlink{0000-0000-0000-0000}~[{PUC Minas}~|\href{mailto: joao.estevao@sga.pucminas.br}{~{\textit{joao.estevao@sga.pucminas.br}}}~]}
\affil{\textbf{Letícia Azevedo Cota Barbosa}~\orcidlink{0000-0000-0000-0000}~[{PUC Minas}~|\href{mailto: lacbarbosa@sga.pucminas.br}{~{\textit{lacbarbosa@sga.pucminas.br}}}~]}
\affil{\textbf{Lívia Alves Ferreira}~\orcidlink{0000-0000-0000-0000}~[{PUC Minas}~|\href{mailto: livia.alves@sga.pucminas.br}{~{\textit{livia.alves@sga.pucminas.br}}}~]}
\affil{\textbf{Luiz Fernando Antunes da Silva Frassi}~\orcidlink{0000-0000-0000-0000}~[{PUC Minas}~|\href{mailto: lfasfrassi@sga.pucminas.br}{~{\textit{lfasfrassi@sga.pucminas.br}}}~]}
\affil{\textbf{Pedro Augusto Gomes Ferreira de Albuquerque}~\orcidlink{0000-0000-0000-0000}~[{PUC Minas}~|\href{mailto: pagfalbuquerque@sga.pucminas.br}{~{\textit{pagfalbuquerque@sga.pucminas.br}}}~]}
\affil{\textbf{Pedro Henrique Lopes Costa}~\orcidlink{0000-0000-0000-0000}~[{PUC Minas}~|\href{mailto: pedro.costa.1362465@sga.pucminas.br}{~{\textit{pedro.costa.1362465@sga.pucminas.br}}}~]}
\affil{\textbf{Victor Souza Lima}~\orcidlink{0000-0000-0000-0000}~[{PUC Minas}~|\href{mailto: victor.lima.1478515@sga.pucminas.br}{~{\textit{victor.lima.1478515@sga.pucminas.br}}}~]}
}

\begin{document}

\begin{frontmatter}

\maketitle

\begin{mail}
Pontifícia Universidade Católica de Minas Gerais, Av. Dom José Gaspar, 500 - Coração Eucarístico, Belo Horizonte, Minas Gerais, 30535-901, Brasil.
\end{mail}

\begin{abstract-pt}
Este trabalho relata a segunda entrega do Voxel Royale, um mini \textit{battle royale} voxel com cliente web e backend distribuído. A versão atual mantém Gateway, Lobby e Game como serviços Go independentes, comunicando-se internamente por gRPC e expondo HTTP/JSON e WebSocket pelo Gateway. Em relação à entrega anterior, foram adicionados cliente Phaser/Vite, partida em tempo real com relógio autoritativo no servidor, integração Lobby $\rightarrow$ Game por \texttt{StartMatch}, salas com jogadores prontos e início automático, observabilidade com Prometheus/Grafana/Jaeger, teste de carga com 50 jogadores e tolerância a falhas no Lobby por replicação primário-backup. O Lobby primário replica snapshots versionados para o backup e só confirma escritas após a réplica aceitar o estado; em falha do primário, o Gateway promove o backup e repete a operação. Testes unitários, smoke tests via HTTP/WebSocket e o runner \texttt{stress50} validaram os fluxos principais, rollback de falhas e 50 conexões simultâneas.
\end{abstract-pt}

\begin{abstract-en}
This paper reports the second delivery of Voxel Royale, a voxel mini battle royale with a web client and a distributed backend. The current version keeps Gateway, Lobby and Game as independent Go services, using gRPC internally and exposing HTTP/JSON and WebSocket through the Gateway. Compared with the previous delivery, the project now includes a Phaser/Vite client, real-time gameplay driven by an authoritative server clock, Lobby-to-Game integration through \texttt{StartMatch}, ready-state rooms with automatic start, Prometheus/Grafana/Jaeger observability, a 50-player load runner and primary-backup fault tolerance for the Lobby. The primary Lobby replicates versioned snapshots to the backup and only confirms writes after the replica accepts the state; when the primary fails, the Gateway promotes the backup and retries the operation. Unit tests, HTTP/WebSocket smoke tests and the \texttt{stress50} runner validated the main flows, failure rollbacks and 50 simultaneous connections.
\end{abstract-en}

\begin{pchaves}
Computação distribuída, tolerância a falhas, replicação primário-backup, microsserviços, gRPC, WebSocket, jogos multiplayer.
\end{pchaves}

\begin{keywords}
Distributed computing, fault tolerance, primary-backup replication, microservices, gRPC, WebSocket, multiplayer games.
\end{keywords}

\end{frontmatter}

\section{Introdução}
\label{sec:intro}

Jogos multiplayer em tempo real precisam combinar baixa latência, consistência do estado compartilhado e resiliência a falhas parciais. No \textit{Voxel Royale Distribuído}, esses problemas aparecem em um cenário controlado: até 50 jogadores entram em uma sala, iniciam uma partida curta e recebem snapshots autoritativos de uma arena com movimento, baús, armas, dano, zona segura e ranking.

A entrega anterior demonstrava a divisão inicial em microsserviços, gRPC~\cite{grpc}, \textit{Protocol Buffers}~\cite{protobuf}, Web Services HTTP via \texttt{grpc-gateway}~\cite{grpcgateway} e empacotamento com Docker Compose~\cite{dockercompose}. A segunda entrega preserva essas seções e amplia a implementação para cobrir requisitos de tempo real, teste de carga, observabilidade e tolerância a falhas. O sistema completo agora sobe com frontend, Gateway, Game, Lobby primário, Lobby backup, Prometheus, Grafana e Jaeger.

O foco deste texto é registrar as novas implementações que ainda não estavam no artigo: a replicação primário-backup do Lobby, o failover feito pelo Gateway, o tratamento de falhas nas transições Lobby $\rightarrow$ Game, o pipeline WebSocket $\rightarrow$ gRPC em tempo real, os testes executados e as melhorias futuras.

Como contribuição prática, o projeto demonstra uma arquitetura distribuída pequena, mas executável de ponta a ponta. A aplicação permite observar os efeitos de decisões clássicas de sistemas distribuídos em um domínio interativo: a consistência do estado de coordenação fica no Lobby, a autoridade temporal da simulação fica no Game, e a borda pública concentra adaptação de protocolos, monitoramento e recuperação diante da indisponibilidade do serviço primário de salas.

\section{Arquitetura do Sistema}

O backend permanece organizado em três serviços principais, todos escritos em Go e executados em containers separados:

\begin{itemize}
\item \textbf{Gateway} (\texttt{:8080}): única borda pública. Expõe HTTP/JSON, \texttt{/healthz}, \texttt{/readyz}, \texttt{/metrics} e o WebSocket \texttt{/v1/match/ws}. Encaminha chamadas para Game e Lobby por gRPC.
\item \textbf{Lobby} (\texttt{:50052} gRPC, \texttt{:8081} HTTP interno): gerencia salas, jogadores, dono da sala, marcação de pronto, início manual ou automático e replicação do estado de salas.
\item \textbf{Game} (\texttt{:50051} gRPC, \texttt{:8082} HTTP interno): mantém partidas por \texttt{room\_id}, executa o relógio autoritativo de 15 Hz e publica snapshots de \texttt{GameState}.
\end{itemize}

A topologia da entrega atual adiciona um \texttt{lobby-primary} e um \texttt{lobby-backup}. O Gateway conversa primeiro com \texttt{lobby-primary:50052}; o primário replica para \texttt{http://lobby-backup:8081/replication/lobby-state}. Se o primário fica indisponível, o Gateway chama \texttt{POST /replication/promote} no backup, troca o endpoint ativo para \texttt{lobby-backup:50052} e repete uma vez a operação que falhou.

O fluxo distribuído principal possui quatro etapas. Primeiro, o cliente cria ou entra em uma sala por HTTP no Gateway; a chamada é traduzida para gRPC e persistida em memória pelo Lobby. Segundo, os jogadores marcam \texttt{ready}; quando todos estão prontos, o Lobby chama \texttt{StartMatch} no Game com o \textit{roster} da sala. Terceiro, cada navegador abre um WebSocket no Gateway, que envia inputs para \texttt{PushInput} e recebe snapshots por \texttt{WatchMatch}. Por fim, métricas e traces registram a passagem por Gateway, Lobby e Game, permitindo verificar latência, erros e volume de mensagens durante a demonstração.

\subsection{Gateway Service}

O Gateway mantém o caminho HTTP $\rightarrow$ gRPC da primeira entrega, mas o fluxo do Lobby deixou de ser apenas o \texttt{grpc-gateway} gerado. Para suportar failover, \texttt{internal/gateway/lobby\_failover.go} implementa um proxy HTTP que decodifica as rotas de sala, chama o cliente gRPC do Lobby ativo e converte erros gRPC em status HTTP. O failover só é disparado para falhas de disponibilidade, especificamente \texttt{Unavailable}, \texttt{DeadlineExceeded} ou timeout de contexto.

O tempo real usa um handler WebSocket dedicado. O navegador abre \texttt{GET /v1/match/ws?room=\{id\}\&player=\{id\}}; o Gateway reescreve \texttt{room\_id} e \texttt{player\_id} a partir da conexão, ignora esses campos vindos do payload do cliente e encaminha cada input por \texttt{Game.PushInput}. Em paralelo, assina \texttt{Game.WatchMatch} e envia os snapshots recebidos de volta ao WebSocket em JSON. Pings periódicos e \textit{deadlines} de leitura/escrita detectam conexões quebradas sem travar o processo.

\subsection{Lobby Service e Replicação}

O Lobby continua mantendo o ciclo de vida das salas em memória: \texttt{CreateRoom}, \texttt{JoinRoom}, \texttt{GetRoom}, \texttt{StartRoom}, \texttt{LeaveRoom} e \texttt{SetReady}. A novidade funcional é que \texttt{SetReady} marca cada jogador como pronto e inicia automaticamente a sala quando todos estão prontos. Tanto o início manual quanto o automático chamam \texttt{Game.StartMatch}, passando o \textit{roster} da sala para pré-posicionar os jogadores no Game.

A replicação implementada é primário-backup. O papel de cada instância vem de \texttt{LOBBY\_REPLICATION\_ROLE}; o primário recebe escritas públicas e o backup é somente leitura até ser promovido. A cada escrita bem-sucedida (\texttt{CreateRoom}, \texttt{JoinRoom}, \texttt{LeaveRoom}, \texttt{SetReady} ou \texttt{StartRoom}), o primário incrementa uma versão monotônica, gera um snapshot completo com salas, jogadores, dono, status, flags de pronto e contadores de ID, e envia esse evento ao backup por HTTP interno.

O backup aceita apenas \texttt{version == lastVersion + 1}. Eventos duplicados com a mesma versão são idempotentes somente quando o snapshot é igual ao estado já aplicado; eventos fora de ordem ou com lacuna são rejeitados com erro. O primário só responde sucesso ao cliente depois da confirmação do backup. Se a replicação falhar, \texttt{commitReplicationLocked} restaura o snapshot e a versão anteriores, retornando \texttt{Unavailable}. Assim, uma sala confirmada pelo Lobby também está presente no backup.

\subsection{Game Service e Cliente}

O Game agora mantém um mapa de partidas por \texttt{room\_id}, além de uma partida global de compatibilidade com a rota unária \texttt{/v1/match/stream}. \texttt{StartMatch} cria ou reinicia a partida de uma sala com os jogadores recebidos do Lobby; \texttt{PushInput} armazena o último input por jogador; \texttt{WatchMatch} envia um stream de snapshots. O relógio da partida começa quando o primeiro assinante entra e para quando o último sai, evitando CPU ociosa em partidas sem clientes.

Cada tick aplica inputs ainda válidos, descarta sequências antigas, limita movimento, evita colisão com pedras, processa abertura de baús, ataque, dano de zona segura e ranking. O jogo dura até 4500 ticks, cerca de cinco minutos a 15 Hz, ou termina quando resta no máximo um jogador vivo. O cliente web em Phaser/Vite consome esse fluxo via WebSocket, renderiza lobby, sala, arena e tela de fim, e possui fallback local para demonstração quando o backend não responde.

\subsection{Infraestrutura e Observabilidade}

O Docker Compose atual sobe frontend, Gateway, Game, Lobby primário, Lobby backup, Prometheus, Grafana e Jaeger. \texttt{/readyz} diferencia processo vivo de serviço pronto: o Gateway só fica pronto quando consegue alcançar o Game e pelo menos um Lobby utilizável; os Lobbies verificam a dependência com o Game; o Game fica pronto após abrir o listener gRPC.

As métricas Prometheus~\cite{prometheus} incluem sessões WebSocket, mensagens e bytes por direção, latência de \texttt{PushInput}, erros realtime, tentativas de failover, inputs aceitos/rejeitados, streams ativos, ticks, duração do tick, snapshots descartados, salas, jogadores e eventos de replicação. O Jaeger~\cite{jaeger} recebe traces OpenTelemetry~\cite{opentelemetry} das entradas HTTP e chamadas gRPC.

\subsection{Decisões de Projeto}

A separação entre Lobby e Game reduz acoplamento entre coordenação e simulação. O Lobby mantém dados de controle de baixa frequência, como dono, lista de jogadores, capacidade e estado de pronto; por isso é viável replicar snapshots completos a cada escrita. O Game, por outro lado, produz estado em alta frequência e aceita perda controlada de snapshots para clientes lentos, desde que o próximo estado autoritativo substitua o anterior.

O Gateway foi escolhido como ponto único de adaptação porque simplifica o cliente e preserva contratos internos tipados. Externamente, o sistema oferece HTTP/JSON para operações de sala e WebSocket para tempo real. Internamente, os serviços usam gRPC, o que torna explícitos os campos de entrada e saída, facilita geração de clientes e mantém a comunicação entre containers independente da tecnologia do frontend.

O modelo primário-backup adotado prioriza consistência simples em vez de disponibilidade máxima. Durante operação normal, uma escrita no Lobby só é confirmada após o backup aplicar o snapshot correspondente. Após promoção, o backup passa a aceitar escritas sem uma segunda réplica, o que mantém o sistema utilizável para a demonstração, mas deixa claro que a tolerância cobre a falha do primário, não múltiplas falhas simultâneas.

\section{Implementação dos Requisitos}

\subsection{Comunicação, Web Services e Tempo Real}

Os contratos vivem em \texttt{proto/lobby/v1/lobby.proto} e \texttt{proto/match/v1/match.proto}; o código Go gerado fica em \texttt{gen/}. As rotas HTTP públicas do Lobby são \texttt{POST /v1/rooms}, \texttt{POST /v1/rooms/\{id\}/join}, \texttt{GET /v1/rooms/\{id\}}, \texttt{POST /v1/rooms/\{id\}/start}, \texttt{POST /v1/rooms/\{id\}/leave} e \texttt{POST /v1/rooms/\{id\}/ready}. O Game expõe \texttt{POST /v1/match/stream} para compatibilidade e usa internamente \texttt{StartMatch}, \texttt{PushInput} e \texttt{WatchMatch} para a partida em tempo real.

Essa separação mantém HTTP/JSON como interface pública de operação e gRPC como transporte interno tipado. A partida em tempo real não depende de polling HTTP: o WebSocket do Gateway representa a sessão do jogador, enquanto o Game concentra a autoridade do estado e publica snapshots regulares.

\subsection{Modelagem das Mensagens}

O contrato do Lobby representa salas como entidades de coordenação. Cada resposta de sala contém \texttt{room\_id}, \texttt{status}, \texttt{owner\_id}, \texttt{max\_players}, \texttt{join\_url} e a lista de jogadores com identificador, nome e estado de pronto. Essa modelagem permite que o frontend renderize a sala antes da partida e que o Gateway repita operações após failover sem depender de estado local do navegador.

No contrato do Game, \texttt{PlayerInput} carrega movimento, sequência, intenção de abrir baú e ataque. O campo \texttt{input\_sequence} permite descartar comandos antigos que chegam fora de ordem, enquanto \texttt{room\_id} e \texttt{player\_id} isolam as partidas e impedem que inputs de uma sala afetem outra. O snapshot \texttt{GameState} agrega \texttt{tick}, jogadores, baús, zona segura, ranking e ticks restantes, formando a visão autoritativa que o cliente apenas renderiza.

O uso de \textit{Protocol Buffers} foi importante para manter essa evolução controlada. Novos campos puderam ser adicionados aos contratos sem reescrever manualmente serialização, clientes gRPC ou o proxy HTTP gerado. Isso reduziu o risco de divergência entre documentação, backend e frontend durante a segunda entrega.

\subsection{Mapeamento aos Requisitos da Entrega}

A Tabela~\ref{tab:requisitos} resume como a implementação atende aos principais pontos avaliados na segunda entrega. A evidência combina código, execução local via Docker Compose, testes automatizados e resultados documentados de carga.

\begin{table}[ht]
\centering
\caption{Relação entre requisitos distribuídos e implementação.}
\label{tab:requisitos}
\scriptsize
\begin{tabularx}{\columnwidth}{|X|X|}
\hline
\textbf{Requisito} & \textbf{Atendimento no Voxel Royale} \\ \hline
Serviços independentes & Gateway, Lobby e Game possuem entrypoints, Dockerfiles, portas e responsabilidades próprias. \\ \hline
Comunicação distribuída & HTTP/JSON e WebSocket na borda; gRPC e \textit{Protocol Buffers} entre serviços. \\ \hline
Tratamento de falhas & Health/readiness, rollback de \texttt{StartRoom}, replicação primário-backup e promoção do Lobby backup. \\ \hline
Tempo real & Relógio autoritativo no Game a 15 Hz, \texttt{PushInput}, \texttt{WatchMatch} e WebSocket no Gateway. \\ \hline
Observabilidade & Métricas Prometheus, dashboard Grafana e traces OpenTelemetry enviados ao Jaeger. \\ \hline
Teste de carga & Runner \texttt{stress50} com 50 WebSockets simultâneos e coleta de inputs, snapshots, bytes e erros. \\ \hline
\end{tabularx}
\end{table}

\subsection{Tratamento de Falhas Implementado}

Foram implementados quatro mecanismos principais. Primeiro, os \textit{healthchecks} e \texttt{/readyz} impedem que o Compose e o Gateway considerem pronto um serviço que ainda não abriu suas dependências. Segundo, o Lobby reverte a transição de sala para \texttt{STARTED} quando \texttt{Game.StartMatch} falha, deixando a sala novamente em \texttt{WAITING}. Terceiro, a replicação primário-backup confirma escritas somente após a réplica aceitar o snapshot, com rollback local se o backup falhar. Quarto, o Gateway promove o backup quando o primário falha por indisponibilidade e mantém a decisão de forma aderente, sem voltar automaticamente ao primário antigo, reduzindo o risco de \textit{split-brain}.

Também há tratamentos no plano de jogo: inputs sem \texttt{player\_id} ou com \texttt{NaN}/\texttt{Inf} são rejeitados; \texttt{input\_sequence} antigo é ignorado; cliente WebSocket lento tem snapshots descartados quando o buffer enche, evitando bloquear o relógio global; desconexões removem o input pendente e encerram o relógio quando não há assinantes.

Essas escolhas delimitam claramente o modelo de falhas assumido. O sistema tolera queda do Lobby primário depois que seu estado foi replicado, indisponibilidade temporária do Game durante a tentativa de início de sala e desconexões de clientes individuais. Não há, nesta entrega, garantia de recuperação de uma partida ativa se o Game cair, nem consenso distribuído para múltiplos candidatos a primário. Essas limitações foram mantidas explícitas para que a implementação continue compatível com o escopo acadêmico e demonstrável em uma máquina local.

\subsection{Testes Realizados}

A suíte \texttt{go test ./...} cobre os três serviços. No Lobby, os testes validam criação/entrada/saída de sala, capacidade, autorização do dono, transferência de posse, \texttt{SetReady}, início automático, chamada a \texttt{StartMatch}, rollback quando o Game falha, backup somente leitura, aplicação de evento de replicação, rejeição de evento fora de ordem e rollback quando a réplica não confirma. No Gateway, há teste de promoção do backup com primário indisponível, verificando que a sala é criada no backup promovido e que o endpoint ativo muda para ele. No Game, os testes cobrem snapshot autoritativo, pré-spawn do roster, isolamento por \texttt{room\_id}, reinício de partida, movimento limitado, colisão com pedras, baús, dano, eliminação, ranking, zona segura, relógio realtime, descarte de input antigo, parada do relógio sem assinantes e limpeza de conexão.

Além dos unitários, foram executados smoke tests por HTTP para \texttt{/healthz}, \texttt{/readyz}, criação/entrada/início de sala e \texttt{/v1/match/stream}. O runner \texttt{tools/stress50} cria uma sala, adiciona bots, inicia a partida e abre 50 WebSockets. O smoke local registrado em \texttt{docs/stress-results.md} executou 50 jogadores por 5 segundos, com 3750 inputs enviados, 3750 snapshots recebidos, média de 75 snapshots por jogador, aproximadamente 15 Hz e zero erros do simulador.

Na validação operacional, o Compose foi usado como ambiente de integração porque exercita DNS interno, healthchecks e dependências entre containers. O cenário de failover foi validado parando o \texttt{lobby-primary} e repetindo uma operação de sala pelo Gateway; o resultado esperado é a promoção do backup, a repetição da chamada e o incremento da métrica \texttt{voxel\_gateway\_lobby\_failovers\_total}. Para observabilidade, Prometheus coleta métricas em \texttt{:8080}, \texttt{:8081} e \texttt{:8082}, enquanto o Jaeger permite inspecionar spans dos serviços \texttt{voxel-gateway}, \texttt{voxel-lobby} e \texttt{voxel-game}.

\subsection{Cenários de Demonstração}

A demonstração recomendada começa pela execução completa com \texttt{make docker-up}, que aciona o Compose da pasta \texttt{deployments/}. Em seguida, o avaliador pode abrir o frontend em \texttt{http://localhost:5173}, criar uma sala, compartilhar a URL ou QR Code, marcar jogadores como prontos e observar a transição automática para a partida quando todos estiverem prontos.

O segundo cenário exercita a comunicação em tempo real. Durante a partida, cada cliente envia movimento e ações pelo WebSocket do Gateway, enquanto o Game aplica esses inputs no relógio do servidor e devolve snapshots. O frontend mostra movimento, zona segura, baús, dano e ranking a partir do estado recebido, sem calcular localmente a verdade da simulação.

O terceiro cenário provoca falha controlada do Lobby primário. Ao parar o container \texttt{lobby-primary} e criar ou consultar uma sala pelo Gateway, espera-se que o backup seja promovido e que as próximas chamadas passem a usar \texttt{lobby-backup:50052}. Esse teste evidencia a diferença entre falha de processo e perda de estado: como as escritas confirmadas já estavam replicadas, o backup possui dados suficientes para continuar o fluxo de salas.

Por fim, o cenário de carga usa \texttt{make stress50}, com parâmetros ajustáveis no runner \texttt{tools/stress50}. O resultado deve ser analisado junto com Prometheus/Grafana para verificar taxa de snapshots, erros, bytes transmitidos e comportamento do relógio do Game sob múltiplas conexões simultâneas.

\section{Desafios Encontrados}

\subsection{Integração Entre Sala e Partida}

Na primeira entrega, \texttt{StartRoom} apenas mudava o status da sala. A integração real exigiu uma interface \texttt{MatchStarter}, um cliente gRPC no processo Lobby e rollback quando o Game não confirma. A chamada foi colocada fora do mutex principal do Lobby para não segurar o lock durante rede, mas a transição de status e sua replicação continuam protegidas por snapshot.

\subsection{Replicação sem Persistência Durável}

Como o escopo não incluía banco de dados, a consistência possível foi concentrada em memória replicada. O desafio foi evitar divergência silenciosa: por isso o backup rejeita lacunas de versão, duplicatas diferentes e escritas públicas enquanto está no papel de backup. A decisão de exigir confirmação síncrona da réplica aumenta a latência de escrita, mas dá uma propriedade simples de demonstrar: se o cliente recebeu sucesso, o backup possui o mesmo estado.

\subsection{Tempo Real e Clientes Lentos}

O fluxo unário da primeira entrega era suficiente para \textit{curl}, mas não para uma partida fluida. A solução foi separar input e snapshot: \texttt{PushInput} atualiza um buffer e o relógio de servidor aplica o último input conhecido a cada tick. Para não deixar um cliente lento bloquear todos os outros, cada assinante possui buffer limitado; em caso de lotação, o snapshot novo é descartado e a métrica \texttt{voxel\_game\_snapshot\_drops\_total} registra a perda.

\subsection{Coordenação de IDs e Estado de Sala}

Outro ponto sensível foi manter IDs estáveis entre criação de sala, entrada de jogadores, início da partida e conexão WebSocket. O Lobby gera identificadores de sala e jogador, replica também os contadores internos e envia o mesmo \textit{roster} para o Game. No Gateway, os parâmetros da URL do WebSocket são tratados como fonte de verdade para \texttt{room\_id} e \texttt{player\_id}, evitando que um cliente envie inputs para outro jogador apenas alterando o payload JSON.

\section{Participação dos Integrantes}

O grupo organizou a implementação por serviço e complementou a divisão com frentes funcionais do jogo, mantendo integração conjunta nos contratos \texttt{.proto}, Docker Compose, testes e tratamento de erros.

\begin{table}[ht]
\centering
\caption{Divisão de responsabilidades por serviço e frente de implementação.}
\label{tab:equipes}
\scriptsize
\begin{tabularx}{\columnwidth}{|X|X|}
\hline
\textbf{Integrantes} & \textbf{Responsabilidade} \\ \hline
Ana Cristina, Lívia e Luiz &
Gateway HTTP, integração HTTP $\leftrightarrow$ gRPC via \texttt{grpc-gateway}, registro dos dois handlers, \textit{healthcheck} e contribuições no personagem, mapa, projéteis, armas e \textit{safe zone} do frontend. \\ \hline
Letícia, Pedro Augusto e Pedro Henrique &
Lobby Service: anotação HTTP do \texttt{.proto}, ciclo de vida de sala em memória (criar, entrar, consultar, iniciar, sair), validações de propriedade/capacidade, implementação do personagem, mapa, sistema de projéteis, design das armas e \textit{safe zone} no frontend. \\ \hline
André, João e Victor &
Game Service: contrato de \texttt{StreamMatch}, regras de movimento, baús, três armas com perfis distintos, \textit{safe zone} progressiva, ranking, integração com o frontend e algoritmo de replicação primário-backup. \\ \hline
Todos os integrantes &
Reestruturação do monorepo, Docker Compose, smoke tests, documentação, observabilidade, revisão do relatório e tratamento de erros. \\ \hline
\end{tabularx}
\end{table}

\section{Conclusão}

A segunda entrega transforma o Voxel Royale de um backend distribuído inicial em uma demonstração completa de comunicação, tempo real, observabilidade e tolerância a falhas. O sistema executa cliente web, Gateway, Game, Lobby primário, Lobby backup e stack de métricas/traces em Docker Compose. O Lobby agora inicia partidas no Game, replica seu estado para um backup versionado e permite failover promovido pelo Gateway. O Game evoluiu para partidas por sala com relógio autoritativo e snapshots via stream gRPC, repassados ao navegador por WebSocket. Os testes unitários, smoke tests e o runner de 50 jogadores validam tanto a lógica funcional quanto os principais cenários de falha implementados.

As principais limitações restantes são deliberadas: o estado do Game ainda não é replicado, não há persistência durável de salas, o failover é manualmente promovido pelo Gateway mas não possui eleição distribuída, e a segurança da borda ainda é simples para fins acadêmicos. Como melhorias futuras, o grupo propõe persistir ou replicar partidas ativas, adicionar logs estruturados com \texttt{request\_id}, \texttt{room\_id} e \texttt{player\_id}, endurecer WebSocket com autenticação leve e controle de origem, tratar reconexão de jogador com mais robustez, rodar testes de carga mais longos em VPS e avaliar uma estratégia de consenso ou lease para promoção automática sem risco de \textit{split-brain}.

Portanto, a entrega atende ao objetivo central de construir e medir um sistema distribuído funcional, com fronteiras de serviço claras e falhas tratadas de forma observável. O resultado não é apenas uma API segmentada em containers, mas uma aplicação interativa na qual o usuário consegue criar salas, iniciar partidas, jogar em tempo real, provocar falhas controladas e acompanhar o comportamento do sistema por métricas e traces.

\begin{declarations}

\begin{materials}
O código-fonte, contratos \texttt{.proto}, Dockerfiles, \texttt{docker-compose.yml}, documentação de replicação (\texttt{docs/lobby-replication.md}), observabilidade (\texttt{docs/observability.md}) e resultados de carga (\texttt{docs/stress-results.md}) estão disponíveis no repositório GitHub: \url{https://github.com/joaodest/comp-dist.git}
\end{materials}

\end{declarations}

\begin{thebibliography}{9}
\bibitem{grpc} gRPC Authors. \textit{gRPC: A high performance, open source universal RPC framework}. \url{https://grpc.io}.
\bibitem{grpcgateway} gRPC-Gateway Authors. \textit{grpc-gateway: gRPC to JSON proxy generator following google.api.http}. \url{https://github.com/grpc-ecosystem/grpc-gateway}.
\bibitem{protobuf} Google. \textit{Protocol Buffers}. \url{https://protobuf.dev}.
\bibitem{dockercompose} Docker Inc. \textit{Docker Compose}. \url{https://docs.docker.com/compose/}.
\bibitem{prometheus} Prometheus Authors. \textit{Prometheus: Monitoring system and time series database}. \url{https://prometheus.io}.
\bibitem{opentelemetry} OpenTelemetry Authors. \textit{OpenTelemetry Documentation}. \url{https://opentelemetry.io/docs/}.
\bibitem{jaeger} Jaeger Authors. \textit{Jaeger: Distributed tracing platform}. \url{https://www.jaegertracing.io}.
\end{thebibliography}

\end{document}
